\documentclass[12pt]{article}  
\newcommand{\say}[1]{``#1''}  
\newcommand{\nsay}[1]{`#1'}  
\usepackage{endnotes}  
\newcommand{\B}{\backslash{}}  
\renewcommand{\,}{\textsuperscript{,}}  
\usepackage{setspace}   
\usepackage{tipa}  
\usepackage{hyperref}  
\begin{document}  
\doublespacing  
\section{\href{daily-reflection-2026-08-03.html}{Just a Daily Reflection}}  
First Published: 2026 August 4\footnote{forgot to post}

N.B. I realized that a 4000 word daily reflection is basically a whole post.
Most of this is entirely contained within the stuff about the embroidery project, which, while still mostly on hold, has made some progress in the realm of design at least.
Woot.


\section{Upcoming Posts}  
\begin{itemize}

\item Wonder  
\item Curse of knowledge, esp re. Music Theory

Most of a draft done, though rambly.  
\item Overview of a course on astrochem

I think that I have writing somewhere about it.  
\item RebelFit

A rambly draft done, and some more work which tells me that I was not doing anything wrong in the past, approximation just takes time.  
\item Why aspiring writers should practice typing\footnote{and this should also like tie in the thing that a lot of writers are told to do which is directly write good poetry or prose down to get it in the hand. This also has the effect of practicing hand writing. I think that there's also the whole like \say{when we want to be better at running, we also stretch }}

\item At some point I'd like to reflect on my hesitance to delete email, which I think comes from a place of digital hoarding. That's going to be a fun one for sure.

\item Moving (homes)

Wrote two partial drafts, one considering the idea of that home being where I last lived that my mother saw.
The second reflected on how I'm finally at the point where radical acceptance feels normal.

\item Algorithmic constraints in art I create\footnote{see footnote in \href{reflection-june-2026}{this post's second daily reflection about written word, poetry}}

\item Singing in a low-voice choir, especially in context of my other music

\item Listening to content at speeds (J speed versus single speed)\footnote{see footnote in \href{reflection-2}{this post's daily reflection about written word, audio book} for context}

\item Book Review(s):\footnote{also added today as a list so I remember}

\begin{itemize}

\item First 20 Hours

\item Dungeon Crawler Carl

\end{itemize}

\item Ideas of Salvation\footnote{from the First Draft of \href{reflection-2}{this post}, footnote where I discuss changes to my theology.}, in particular, how mine have changed over time.

\item Difference between lyrics and poetry in my writing.\footnote{from gut check on writing from \href{daily-reflection-2026-08-03}{this post}}

\item How I make my full to-do list when feeling like there's too much to do.\footnote{from secondary gut check from \href{daily-reflection-2026-08-03}{this post}}

\item Why I favorited the songs I did.\footnote{From how goes listening in \href{daily-reflection-2026-08-03}{this post}}

\end{itemize}



\section{Daily Reflection: 3 August 2026}

\begin{itemize}

\item All in all, how goes it?

Um, I don't think terrible.
I'm finally moved out fully, and more or less\footnote{read: I have everything but the dish drain my lovely partner got me inside my home (and the dishes which are in it)} moved in.
Now I just need to get set up in the new place.
I have already really changed the setup from what I had intended, in large part because my family came to help me move and set one of the bedrooms up as a really nice storage space. I'm still thinking it'll be a good studio workspace, but not exactly how I'd intended.

\item Creative Considerations:

\begin{itemize}

\item Gut and / or vibes check, how is creative life going?

Eh, honestly not great, I haven't been writing much at all, music hasn't been going well, and I haven't embroidered in the least.

\item Written Word:

\begin{itemize}

\item How's your relationship with the written word in general?

Not fantastic.
I have read about 6000 pages of a fan fiction since I started that in um the 22 of July.
That's about 2 weeks, about 3000 a week, about 400 pages a day in average.
Ok, so maybe reading has been going well.

Haven't blogged, haven't written much music, though I did write some lyrics yesterday.\footnote{note to self, blog about difference between lyric and poetry for me.}
Blogging feels like it'll come when I\footnote{finally} set up WiFi in my new apartment, and I have no real answers to the problem with poetry.

\item How goes blogging? What's feeling like the barrier?

It's been nearly a full month.
For a while the issue was feeling overwhelmed and ill.
Now the issue is that I lack WiFi at my home. 

In theory I could go ahead and move to a coffeeshop or something to do the writing\footnote{or, being totally fair, I could do what I'm doing now and not write on the website I like to write}.
I might consider that, especially because there's some work that I need to do these next few days before I get WiFi again.
Otherwise, the issue is really just that I have been feeling overwhelmed.
That's something that blogging and journaling help with, though, so that's something that I need\footnote{need is a strong word, but I'm going to keep it. I want to be my best self, so I do, in fact, need to do this} to be better about requiring for myself.
At the very least, jotting down some thoughts keeps them from spiraling.

\item How goes reading?

As discussed in the header for this section, about 400 pages a day isn't too shabby.
I also have kept up with at least part of the soccer series my brother and I read\footnote{there's the Patreon version that's about a full book ahead of the free releases. Not fully in time with the free releases, because we're about to hit my least favorite part of the series.}, and I've been keeping up with a few others that we both read.
Oh! Shoot, I also read all of another web serial to catch up to where it's now releasing.\footnote{a progression fantasy by an author my brother and I are enjoying}
So, honestly, doing p well there.

\item How goes writing poetry again?

Not great.
As mentioned in the heading question, I wrote some lyric yesterday, and that's great.
Otherwise, I completely spaced that this was something I was doing.

\item We're counting audio book here, anything good?

About halfway through the next Dungeon Crawler Carl book, but that's not been progressing in a hot second.

\item Upon reflection, how's the written word? Anything missing in the four questions?

Eh, about what I thought.
Reflecting a little more, I did read a lot more than I thought I had been, and that's kind of nice to realize.
I just wish that I had those last 1000 pages read already so that I could move on and start reading other stuff again.

Something missing from the questions: hand writing/journaling.

I've been given an opportunity to spend some more time in a reflective space, and I've used that time for doodling and some note taking.
In doing so, I've also been reminded that I love my pens and their inks.
I found the two missing pens!\footnote{the one with a blade tip and the calligraphy / stub nosed}
I've enjoyed them again.

But, this does mean that I have been doing more hand writing.
I made my super to-do list again\footnote{note to self: blog about the super to-do list sometime}, reflected on a few of the tasks, and just generally journaled for a minute.
In general, I think that hand journaling is a good thing for me to do, so I'm going to add that question to the template.

\end{itemize}

\item Music

\begin{itemize}

\item In general, how goes music?

Hmm. 

Not creating much at all, but I did jam with a friend somewhat recently and make (admittedly very tentative) plans to do an open mic together.
I've been listening to a lot of it, and I've been thinking about making it again.
I've gotten really good at drawing treble clefs\footnote{I think for me, at least, starting with the center and working out makes them both look better and more like the G that they're supposed to be}

\item How goes recording?

Nonexistent. I haven't even made a pretense of setting up the studio for recording.

\item How goes writing?

As mentioned in the written word section, haven't done much of any, though a little bit of lyric.
I have been wanting to do more, especially choral and especially low-voice writing.
Also, since my alarm is a short four-part arrangement of three-fishers for digital strings, I've been wanting to get back into that a little more.
Might be worth considering, especially since I know that what I've been loving in the more vocally interesting\footnote{read: more vocal harmonies and emphasis on them} that I've been hearing is the chord movement.
I normally don't plot out my music at all like that, but there's a reason we spend so long doing so in music theory.

I wish that I was better at transcription, I think.
Then again, I don't know if it's actually something I enjoy or just something I like being able to see.
I'm sure there's elements where learning to transcribe makes me better at hearing where I've gone wrong, and so I might consider getting back into ear training.
Sure, with all my copious extra energy.

\item How goes playing?

Haven't really been.
Noodled a bit on the dulcimer and kalimba and music maker\footnote{really need to figure out what that instrument is really called} for a bit recently, but not really much at all.
I miss making music, maybe do that tonight?

\item How goes listening?

Honestly really good.
I've been keeping up on album club.
I have listened through the weekly tracks that Apple Music gives me.\footnote{I know that I saw something about how someone doesn't want their music given to them by an algorithm, and I respect that. That being said, though, I do appreciate being able to find artists through the algorithm and then listening to more of their discography.}
I also started listening back through the recently favorited songs I have.
Part of me does really want to make a note here or somewhere about what in each song I find compelling.
Hmmm.

Let's make that today's post!

Eh, let's do that in the future.\footnote{added to the list}

\item Upon reflection, how goes music? Anything missed in those four questions?

Honestly music going well from consuming, not as well for producing.
Could ask about part writing and ear training.
Could ask about the e-recording, though I guess that kind of goes both in writing and recording.
Then again, since I think it goes in both, maybe it does need its own post.

\end{itemize}

\item Embroidery Project?

\begin{itemize}

\item In general, what's the vibe of the project right now?

Honestly, on hold is the general vibe.
I worked it a bit when visiting my partner a few weeks ago, and I have not touched it since.

\item How goes progress on the design? What's in the way?

Design is honestly going pretty well.
I realized a few things:

\begin{itemize}
\item I want gaps in the braid
\item I don't think I'll be able to do a repeating motif, given that there's so many curves
\item I want even spacing along the knot, rather than along a perfect line.
\item There's a way to make a dot appear evenly spaced along a curve in Procreate\footnote{the app I use for drawing and general design}
\item Woo!
\end{itemize}

So, at this point I know what the braid looks like, I just need to figure out how many stitches wide to make them, etc.
I also have a question to look up, because my partner's aunt, who introduced me to embroidery tapestry\footnote{reverse the words in your head} apparently doesn't have all stitches the same direction at all times.
Is that just her or a general thing?

I mean I'm not doing anything majorly time-sensitive right now.
Let's look now.\footnote{thanks cell data}

Well, a quick internet search doesn't seem to provide answers...
I have mixed feelings about doing so, but I don't entirely know what the right answer is.
Maybe message a friend?
Who do I know who might own a tapestry embroidery kit?

Sadly seems like searching is victim of internet generative\footnote{remember, not intelligent} slop.
So, I might have to go to a meeting of the local embroidery guild sometime soon.

Ok having answered the weave layer question\footnote{ooh non linear writing spicy}, let's go to Wikipedia for counted thread embroidery.
Unfortunately, tapestry embroidery doesn't appear as its own page :(.
\begin{itemize}
\item Assisi: outlines of shapes rather than filling them in.
Uses what appears to be a repeating shape for non-borders.
Borders are done in so that the stitches are lying along the border\footnote{so horizontal line is horizontal stitch}

My impression: Stitches can vary in neither direction nor size, except the borders, where one wants similarly sized stitches.
However, doesn't fill in the space.

\item Bargello: Upright and flat.
Also called Florentine work, Hungarian point, Flame stitch, and Irish stitch.
Apparently, strictly speaking, the flame motif is also essential to the upright stitch, and the basic size is four stitches, though since the beginning of the craft they have varied in size, even within the same design.

My impression: All vertical stitches, but variable number of holes skipped. Space filling.

\item Blackwork: absolutely anything goes.
All in black, also called Spanish blackwork.

Note to self: Holbein Stitch could be good when doing long works.

My impression: Anything goes, size or direction can change. Monochrome and doesn't fill in the full space.

\item Kasuti: same as blackwork, but in color.
Indian style of embroidery, linked from Blackwork page instead of counted thread.

My impression: Anything goes. Multicolor enabled, the thread is the design, and so not space-filling.

\item Canvas work/ needlepoint: uhhhhh generally a catch-all I think.
Or, at least, the oldest form and therefore most consuming.
Often just uses a tent stitch\footnote{half a cross-stitch} along full fabric, like I'd thought.
However, many other options including just a bunch of space-filling methods.
The question then becomes how much I want to have the embroidery stand as a thing that varies solely by my color choice and how much I want the thread to do the speaking.
Ughhh this is like the issue with transparency, isn't it?
Adding a note for this.

\end{itemize}

Well, ok so it seems like it's up to me.
I'm going to go with my initial instinct for this project, which has me not varying the lines?
Hmm, gotta make a few test swatches. Questions for the test swatch(es):

\begin{itemize}
\item How do I cross lines? 

That is, how do I fill space when one set of stitches goes front slash and another goes backslash?
\item Does it work to have more than just the two directions of front and back at 45 degrees? 

That is, can I have the knot design and braid have stitches which move more or less completely along their direction of movement, rather than the thread in place?
\item If I have more than two directions, how does that work to fill space?
\item Do I want the weave to be in more than one direction?

That is, in general do I want the two diagonals to cross each other and have front/back slashes or do I want to deal with all front slashes like I had planned to this point?

Also, if I want more than just front slashes, do I want to add breaks like I'd considered, and therefore also add in potentially more directions with that?
\item Do I want the base layer to be in more than one direction?

In general, tent stitch or some variation would seem the most sensible.
However, there's also something to be said about the idea of the directionality of the weave or the knot impacting the canvas below.
I think that I'm going to lean towards base layer is just in the one direction, at least as much as is viable given the whole \say{how do I cross lines?} question.
\end{itemize}

These'll be fun questions to draw and then stitch out!

Ope, may need to re-answer the stitch questions again.
I don't think that I'll need a test swatch for each question, but I do think that I will need multiple test swatches.

\item How goes progress on the stitching? What's in the way?

At this point, just lacking the design to do the stitching.
When I have that, I'll do some test swatches again.

Update after finishing the design reflection: I think that I need to make some more test swatches.
Need to draw up how they'll look, and need to find the embroidery canvas that I was using.\footnote{or get some more}

\item Upon reflection, how goes embroidery? Anything missed in those two questions?

On hold generally.
That's fine, though.
I'll get back to it when it's the time to do so.

\item Project Progress

\begin{itemize}

\item Base layer design:

Still in ideation mode.

No update.

\item Weave layer design:

Have demonstrated crossing with three stitches wide and four stitch gap at widest point. Looks nice, shows enough of the background that I'm happy with it, probably.  
Will need to consider in context of the top layer,

Also will need to consider in the context of potentially not having all of the stitches going the same direction.
I get why, when doing a less geometric design in particular, one might not be as set on simply using the same stitch over and over.
I just don't know how that works with counted thread embroidery...

Let's return to the question above about the design.

Ok so lotta open questions there. Mainly just about when if ever to add breaks, and whether to use front and back slashes or just front slashes.

\item Top layer design:

Have the knot made\footnote{thanks isometric drawing guide!}, have the spacing for the braid made\footnote{thanks ability to add spacing between points in a pen}.
Even have a vague approximation of what the braid will look like over one third of the project.
Still need:

\begin{itemize}
\item Figure out if i want to have multiple directions in the stitches.
\item Figure out the number of stitches wide to generally aim for for the strands of the braid.  
\item Figure out how to have the fact that the braids are two-colors thick work with the braid turning.\footnote{In general, this is an easy enough problem, I can draw what the ribbon would look like. The issue is pixelization and making that look right on canvas and also not be infinitely many stitches}
\item Fill in the braid for the first third of the project.
\end{itemize}

\item Full design: Figure out how many colors of thread I'm going to buy and how much of each
\begin{itemize}
\item Need six colors for the knot
\item I think that I want to have one color for the weave, though that may end up being one mixed-color color.\footnote{e.g. Gold and green}
\item Figure out how many colors for the background.\footnote{Still leaning towards some sort of gradient but need to make it work with the rainbow.}
\end{itemize}

\item Full design: Figure out how it all looks together. 

That is, take the full stitch design for the knot and add the stitches for the weave and background underneath.

\item Full design: Figure out the action plan for doing the stitches.

Since every stitch is going to be marked out\footnote{or, at least, in general every stitch should be mostly mapped, though maybe not all for the weave and certainly not all for the background}, I'm thinking just go one color at a time?
Like for the knot layer, do all the red, then all the orange, and so on.
That'll look cool, to be certain!

Or, do I start with the background and work my way to solid color?
Eh, it'll end up depending how the swatches come out for multi-lines.
If it ends up that all I want is the single forward slash tent stitch, that'll make it much easier to do the background first.
Eh, that's questions for future.

\item Stitching: do.

I'll be curious what the rate that I can stitch ends up being. I'll have an idea of how many stitches large the project is, so then we can do time per stitch times n stitches is total time.
I'm guessing that it'll end up being about a stitch every ten seconds, especially once we account for threading the needle. 

Then again, I'm told that a needle threader exists, and that might save massive amounts of threading time.

I've mentioned it in the past, but I fight the two forces of:
\begin{enumerate}
\item Threading the needle takes a lot of time and so I want the longest possible thread to not have to do it as often. Plus, there's wasted thread from the ends I need to keep. Therefore, maximize working thread
\item Longer working thread takes longer to pull through, tangles and occasionally knots, and so I want a shorter working thread.
\end{enumerate}
Thread is honestly cheap enough that losing bits for the stabler stitches or the ends of shorter working lengths shouldn't really be a bother to me. That being said, we'll see how good a threader is before I make any other decisions.
I'll have to ask my friends how to use one as well, since the directions seem spooky.

\end{itemize}

\item Things to try with embroidery:\footnote{new today}
\begin{itemize}
\item Using variety of stitches for shaping
\item Using colors for transparency or to give the idea of light better
\item Representational (or non-non-representational to be honest to myself for thoughts) work
\item Smaller projects
\item Using whitespace better. Right now I have all designs as much thicker than one stitch.
Lots of the major styles (blackwork and Assisi, for example) base a lot on single stitch design. Why don't we try some of that?
\end{itemize}

\end{itemize}

\item Outreach:

\begin{itemize}

\item How's outreach going?

Nonexistent. I did get an email recently though.

\item How goes prep for the rest of your life?

Need to do more of that.

\item How goes prep / submission of workshops?

... whoop

\item Any talks coming?

Nope!

\item Upon reflection, how goes outreach? Anything missed in those three questions?

I need to do more again.

\end{itemize}

\item Upon reflection, how is creation going?

Generally decently.
Also crocheted a few balls recently, and in doing so I managed to convince some people to buy memory foam scrap.
Also made a flower for a bartender.


\end{itemize}

\item Health Considerations:

\begin{itemize}

\item Gut and / or vibes check, how is health going?

Physical health good, mental health bad.

\item Working out?

Doing great here! Averaging one or two Pilates classes a day.
Double headers are a killer but I leave them feeling great.

\item Food?

Eating more than I had been a few weeks ago, still not where I want to be.

\item Water?

Not really where I want to be in general. Still, drinking more than before and it's always a game of progress.

\item Sleep?

I was having some insomnia, so tried a lil bit of sleep modulation. I think that the insomnia was an acute thing from an infection, and the need for more sleep is actually the chronic bit. Once I finish this reflection, gonna go take a nap.

\item Upon reflection, how is it going?

Eh! Could be doing better, could be doing worse.

\end{itemize}

\item Practical Considerations:

\begin{itemize}

\item Gut and / or vibes check, how are the practical things going?

I forget what these are. So I'm going with good.


\item How goes packing old home?

Done! Get to remove this question from future answers.


\item How could you be working on your LDR better?

She's really busy right now, so I could definitely be doing more to support her in taking care of herself.

\item How goes setup of new home?

Pretty decently! The public-facing space is mostly viewable, and I'm excited for my open house.

\item Upon reflection, how go practical matters?

Decent!

\end{itemize}

\item Upon total reflection, how goes it?

Honestly, better than I had expected.
I think that when I'm catching up on an old web serial, it can be demoralizing to read for days on end only to still have more to read.
This book at least is constantly moving forward, which I do find pretty impressive, especially for a web serial.
No wonder it's so highly regarded.

Anyways, time for a nap while I consider what to actually blog about today.\footnote{though the voice in my head does say to just cut the embroidery section out of the daily reflection and call it today's post... maybe, depends how we're feeling}
Then again, four thousand words of a daily reflection is also something to be impressed with, even if I didn't have other content.

\end{itemize}

\end{document}